\documentclass[10pt,a4paper]{article}

\usepackage[left=1.85cm,right=1.85cm,top=1.45cm,bottom=1.45cm]{geometry}
\usepackage[T1]{fontenc}
\usepackage[utf8]{inputenc}
\usepackage{libertinus}
\usepackage{microtype}
\usepackage{xcolor}
\usepackage{tabularx}
\usepackage{array}
\usepackage{enumitem}
\usepackage[most]{tcolorbox}
\usepackage[hidelinks]{hyperref}
\usepackage{xurl}
\usepackage{parskip}

\definecolor{Navy}{HTML}{17324D}
\definecolor{Teal}{HTML}{247C83}
\definecolor{Ink}{HTML}{27313A}
\definecolor{Muted}{HTML}{66727C}
\definecolor{Pale}{HTML}{F2F6F7}
\definecolor{PaleBlue}{HTML}{EDF4F6}
\definecolor{Rule}{HTML}{D8E2E6}

\hypersetup{
  colorlinks=true,
  urlcolor=Teal,
  linkcolor=Navy
}

\setlength{\parindent}{0pt}
\setlength{\parskip}{0.48em}
\setlength{\emergencystretch}{2em}
\linespread{1.02}
\pagestyle{empty}
\raggedbottom

\newcommand{\smallcapslabel}[1]{%
  {\fontsize{8.4}{10}\selectfont\bfseries\color{Teal}\MakeUppercase{#1}}%
}

\newcommand{\reviewer}[4]{%
\begin{tcolorbox}[
  colback=Pale,
  colframe=Pale,
  boxrule=0pt,
  arc=1.2mm,
  left=3.2mm,
  right=3.2mm,
  top=2.2mm,
  bottom=2.2mm,
  borderline west={2.2pt}{0pt}{Teal},
  before skip=0.55em,
  after skip=0.55em
]
{\bfseries\color{Navy} #1}
\hfill
{\small\href{mailto:#3}{#3}}\\[-0.1em]
{\small\color{Muted} #2}\\[0.3em]
{\small\color{Ink} #4}
\end{tcolorbox}%
}

\begin{document}

% =========================================================
% PAGE 1: COVER LETTER
% =========================================================

\begin{tabularx}{\textwidth}{@{}X>{\raggedleft\arraybackslash}p{4.4cm}@{}}

\smallcapslabel{Submission cover letter}
&
{\small\color{Muted} September 27, 2026}
\\[0.18em]

{\Large\bfseries\color{Navy}
Knowledge-Based Systems}
&
\\

\end{tabularx}

\vspace{0.55em}
{\color{Rule}\rule{\textwidth}{0.7pt}}
\vspace{0.75em}

{\smallcapslabel{Manuscript}}\\[0.18em]

{\fontsize{14.2}{17}\selectfont\bfseries\color{Navy}
Physics as the label for measuring and correcting materials\\
reasoning in multimodal models\par}

\vspace{0.55em}

\begin{tcolorbox}[
  colback=PaleBlue,
  colframe=PaleBlue,
  boxrule=0pt,
  arc=1.2mm,
  left=3mm,
  right=3mm,
  top=1.9mm,
  bottom=1.9mm
]

\begin{tabularx}{\textwidth}
{@{}>{\centering\arraybackslash}X
>{\centering\arraybackslash}X
>{\centering\arraybackslash}X@{}}

{\bfseries\color{Navy} 6}
{\small physics-oracle families}
&
{\bfseries\color{Navy} 9}
{\small contemporary models}
&
{\bfseries\color{Navy} 4,472}
{\small checkable records}

\end{tabularx}

\end{tcolorbox}

\vspace{0.35em}

Dear Editors,

We are pleased to submit our manuscript,
``\textbf{Physics as the label for measuring and correcting materials reasoning in multimodal models},''
for consideration in \textit{Knowledge-Based Systems}.

Vision--language and language models are increasingly used to interpret
diffraction patterns, microscopy images, spectra, and computational materials
data. However, most current evaluations emphasize final-answer accuracy and
provide limited information about whether the underlying scientific reasoning
obeys established physical constraints. We introduce \textbf{MatPCR}, a
physics-grounded framework that operationalizes established physical knowledge
as executable constraints for evaluating multimodal materials reasoning. Six
programmatic oracle families span diffraction geometry, microscopy scale, Raman
peak structure, near-hull stability, computed band-gap class, and net
magnetization.

Across nine contemporary models and 4,472 checkable reasoning records, MatPCR
reveals substantial variation in physical consistency across both models and
constraint families. We further introduce Constraint-Grounded
Self-Verification, which systematically reduces oracle-detected violations in
all four models tested without regressing previously consistent claims.
Importantly, the accompanying accuracy analysis demonstrates that increased
physical consistency does not necessarily imply increased physical correctness,
establishing a quantitative distinction between physically admissible reasoning
and correct scientific inference. A learned verifier reaches an in-distribution
ROC AUC of 0.850, while its limited transfer to held-out physical constraints
further highlights the need for explicit knowledge-grounded evaluation.

We believe this work is particularly well aligned with
\textit{Knowledge-Based Systems} because it develops a knowledge-grounded
artificial-intelligence evaluation methodology rather than a routine application
of an existing model. Established physical rules are encoded as executable
domain knowledge and used to evaluate and improve model reasoning, directly
connecting machine-learning methodology with knowledge representation and
engineering. The framework is evaluated across multiple physical tasks and data
modalities, including visual materials characterization, and therefore also
provides a practical knowledge-based computer-vision and engineering
application. More broadly, MatPCR offers an auditable approach for testing
whether AI-generated scientific reasoning is consistent with domain knowledge,
even when complete answer-level labels are unavailable.

The complete benchmark, provenance records, prompts, and machine-readable
results are publicly available at
\href{https://zenodo.org/records/21829644}
{Zenodo (record 21829644)}.
The oracle implementations, evaluation code, unit tests, and verifier weights
are available at
\href{https://github.com/KurbanIntelligenceLab/matpcr}
{github.com/KurbanIntelligenceLab/matpcr}.

A preprint version of this work is available on arXiv; the present manuscript
contains the finalized analyses and presentation prepared for journal submission.

This manuscript is original, is not under consideration elsewhere, and has been
approved by all authors. The authors declare no competing interests.

Thank you for considering our manuscript for publication in
\textit{Knowledge-Based Systems}.

\vspace{0.25em}

Sincerely,\\[0.18em]

{\bfseries\color{Navy} Mustafa Kurban}\\[-0.05em]
{\small on behalf of all authors}\\[-0.05em]
{\small Ankara University, Ankara, T\"urkiye}\\[-0.05em]
{\small\href{mailto:kurbanm@ankara.edu.tr}
{kurbanm@ankara.edu.tr}}

\vfill

{\footnotesize\color{Muted}
Manuscript:
\textit{Physics as the label for measuring and correcting materials reasoning in multimodal models}\\
Authors: Hasan Kurban, Rasul Khanbayov, and Mustafa Kurban}

% =========================================================
% PAGE 2: SUGGESTED REVIEWERS
% =========================================================

\clearpage

\begin{tabularx}{\textwidth}{@{}X>{\raggedleft\arraybackslash}p{4.4cm}@{}}

\smallcapslabel{Editorial information}
&
{\small\color{Muted} Suggested reviewers}
\\[0.18em]

{\Large\bfseries\color{Navy} Suggested Reviewers}
&
\\

\end{tabularx}

\vspace{0.55em}
{\color{Rule}\rule{\textwidth}{0.7pt}}
\vspace{0.65em}

{\small\color{Ink}
We suggest the following researchers, whose expertise spans scientific machine
learning, materials informatics, artificial-intelligence evaluation,
knowledge-guided modelling, computational materials science, and rigorous
physics-based validation. To the best of our knowledge, none has a conflict of
interest with the authors.}

\reviewer
{Prof. Bryan M. Wong}
{Department of Chemistry, University of California, Riverside, USA}
{bryan.wong@ucr.edu}
{Computational and machine-learning methods for chemical and materials systems;
well positioned to assess the methodological contribution and physical validity
of the knowledge-grounded evaluation framework.}

\reviewer
{Dr. Aaron Gilad Kusne}
{Materials Measurement Science Division, National Institute of Standards and Technology, USA}
{aaron.kusne@nist.gov}
{Expertise in AI measurement and evaluation, physics-informed machine learning,
and autonomous materials research, directly relevant to knowledge-grounded
scientific-AI evaluation and verification.}

\reviewer
{Dr. Kamal Choudhary}
{National Institute of Standards and Technology, USA}
{kamal.choudhary@nist.gov}
{Develops machine-learning models, databases, and multimodal methods for
materials science and can assess both the AI-evaluation methodology and the
materials-data components of the work.}

\reviewer
{Prof. Shyue Ping Ong}
{Department of Materials Science and Engineering, National University of Singapore, Singapore}
{shyue@nus.edu.sg}
{Expertise in materials informatics, first-principles databases, pymatgen, and
the Materials Project ecosystem, directly relevant to the computational
references and domain-knowledge constraints used in this study.}

\reviewer
{Prof. Heather J. Kulik}
{Department of Chemical Engineering, Massachusetts Institute of Technology, USA}
{hjkulik@mit.edu}
{Expertise in data-driven modelling, electronic-structure calculations, and
rigorous computational validation, relevant to assessing the reliability and
scientific grounding of the proposed evaluation framework.}

\vfill

{\footnotesize\color{Muted}
Manuscript:
\textit{Physics as the label for measuring and correcting materials reasoning in multimodal models}\\
Authors: Hasan Kurban, Rasul Khanbayov, and Mustafa Kurban}

\end{document}